% =====================================================================
%  CMPT 981 Project - Proposal
%
%  Template for all project write-ups. Copy this file for the milestone
%  and final report (e.g. milestone.tex, final_report.tex) and adjust
%  the sections. References are shared across all documents in
%  references.bib (same folder).
%
%  Build:  pdflatex proposal && bibtex proposal && pdflatex proposal && pdflatex proposal
%     or:  latexmk -pdf proposal.tex
%  (Overleaf handles this automatically.)
% =====================================================================

\documentclass[11pt]{article}

% --- Page layout & fonts ----------------------------------------------
\usepackage[margin=1in]{geometry}
\usepackage[T1]{fontenc}
\usepackage[utf8]{inputenc}
\usepackage{lmodern}
\usepackage{microtype}

% --- Math ---------------------------------------------------------------
\usepackage{amsmath, amssymb, amsthm}
\usepackage{mathtools}
\usepackage{bm}

% --- Figures & tables ---------------------------------------------------
\usepackage{graphicx}
\usepackage{booktabs}
\usepackage{subcaption}
\usepackage{xcolor}

% --- Algorithms ---------------------------------------------------------
\usepackage{algorithm}
\usepackage{algpseudocode}

% --- References & links (hyperref/cleveref load last) -------------------
\usepackage[numbers,sort&compress]{natbib}
\usepackage[colorlinks=true, linkcolor=blue, citecolor=blue, urlcolor=blue]{hyperref}
\usepackage[capitalize,noabbrev]{cleveref}

% --- Theorem environments -----------------------------------------------
\newtheorem{theorem}{Theorem}
\newtheorem{lemma}[theorem]{Lemma}
\theoremstyle{definition}
\newtheorem{definition}[theorem]{Definition}

% --- Handy macros -------------------------------------------------------
\newcommand{\R}{\mathbb{R}}
\newcommand{\E}{\mathbb{E}}
\DeclareMathOperator*{\argmin}{arg\,min}
\newcommand{\norm}[1]{\left\lVert #1 \right\rVert}

% Inline TODO notes; comment out the first line and uncomment the
% second to hide them all before submitting.
\newcommand{\todo}[1]{\textcolor{red}{[TODO: #1]}}
% \newcommand{\todo}[1]{}

% --- Title --------------------------------------------------------------
\title{Project Title Goes Here\\[0.3em]
  \large CMPT 981 Project Proposal}
\author{
  Author One \quad Author Two \quad Author Three\\
  Simon Fraser University\\
  \texttt{\{author1, author2, author3\}@sfu.ca}
}
\date{\today}

% =====================================================================
\begin{document}
\maketitle

\begin{abstract}
  One paragraph summarizing the problem, why it matters, the proposed
  approach, and how it will be evaluated. \todo{write abstract}
\end{abstract}

% ---------------------------------------------------------------------
\section{Introduction}
\label{sec:intro}

Motivate the problem and state the main question the project answers.
Cite related work like this~\cite{boyd2004convex} or with author names:
\citet{kingma2015adam} proposed Adam. Refer to other sections with
\cref{sec:method}.

% ---------------------------------------------------------------------
\section{Problem Formulation}
\label{sec:problem}

Define the problem formally. For example, given data
$\{(\bm{x}_i, y_i)\}_{i=1}^n$ with $\bm{x}_i \in \R^d$, we solve
\begin{equation}
  \label{eq:objective}
  \min_{\bm{w} \in \R^d} \; f(\bm{w})
  \coloneqq \frac{1}{n} \sum_{i=1}^{n} \ell\bigl(\bm{w}^\top \bm{x}_i, y_i\bigr)
  + \frac{\lambda}{2} \norm{\bm{w}}_2^2 .
\end{equation}
\Cref{eq:objective} is referenced automatically.

\begin{definition}[Smoothness]
  A differentiable function $f$ is $L$-smooth if
  $\norm{\nabla f(\bm{u}) - \nabla f(\bm{v})}_2 \le L \norm{\bm{u} - \bm{v}}_2$
  for all $\bm{u}, \bm{v}$.
\end{definition}

% ---------------------------------------------------------------------
\section{Related Work}
\label{sec:related}

Summarize prior work and how this project differs. \todo{add references}

% ---------------------------------------------------------------------
\section{Proposed Method}
\label{sec:method}

Describe the approach. Pseudocode is shown in \cref{alg:gd}.

\begin{algorithm}[t]
  \caption{Gradient descent (example)}
  \label{alg:gd}
  \begin{algorithmic}[1]
    \Require initial point $\bm{w}_0$, step size $\eta$, iterations $T$
    \For{$t = 0, \dots, T-1$}
      \State $\bm{w}_{t+1} \gets \bm{w}_t - \eta \nabla f(\bm{w}_t)$
    \EndFor
    \State \Return $\bm{w}_T$
  \end{algorithmic}
\end{algorithm}

% ---------------------------------------------------------------------
\section{Experimental Plan}
\label{sec:experiments}

Datasets, baselines, metrics, and expected results. Example table
(\cref{tab:results}) and figure (\cref{fig:example}):

\begin{table}[t]
  \centering
  \caption{Example results table.}
  \label{tab:results}
  \begin{tabular}{lcc}
    \toprule
    Method        & Accuracy (\%) & Time (s) \\
    \midrule
    Baseline      & 85.0          & 12.3     \\
    Ours          & \textbf{88.5} & 10.1     \\
    \bottomrule
  \end{tabular}
\end{table}

\begin{figure}[t]
  \centering
  % Replace the box with: \includegraphics[width=0.6\linewidth]{figures/loss.pdf}
  \fbox{\parbox[c][4cm][c]{0.6\linewidth}{\centering Figure placeholder}}
  \caption{Example figure.}
  \label{fig:example}
\end{figure}

% ---------------------------------------------------------------------
\section{Timeline}
\label{sec:timeline}

\begin{itemize}
  \item \textbf{Week 1--2:} Literature review and problem setup.
  \item \textbf{Week 3--4:} Implement baselines. \emph{(Milestone)}
  \item \textbf{Week 5--7:} Implement and evaluate proposed method.
  \item \textbf{Week 8:} Write final report.
\end{itemize}

% ---------------------------------------------------------------------
\section{Team Contributions}
\label{sec:contributions}

\begin{itemize}
  \item \textbf{Author One:} \todo{responsibilities}
  \item \textbf{Author Two:} \todo{responsibilities}
  \item \textbf{Author Three:} \todo{responsibilities}
\end{itemize}

% ---------------------------------------------------------------------
\bibliographystyle{plainnat}
\bibliography{references}

\end{document}
